% ============================================================
% Come si usano, argomento 18: legge dei grandi numeri e teorema centrale del limite
% ============================================================
\argomento{Legge dei grandi numeri e teorema centrale del limite}{Formule a p.~\pageref*{f:iid}}[Grandi numeri e
  limite centrale]

\begin{scelta}{Come scegliere la formula}
\renewcommand{\arraystretch}{1.25}
\begin{tabular}{@{}P{0.55\linewidth}P{0.39\linewidth}@{}}
  \toprule
  \textbf{Se il testo\ldots} & \textbf{si usa\ldots} \\
  \midrule
  chiede media o varianza della media di $n$ dati & \vv{mediacampione} \\
  chiede quanti dati servono perché la media sbagli di poco con alta probabilità & con certezza Chebyshev \vv{lgn};
    in modo approssimato, e con meno dati, il limite centrale \vv{tcl} \\
  chiede la probabilità che una somma di tante variabili superi un valore & \vv{sommaiid}, poi \vv{tcl} e la
    funzione $Q$ \\
  chiede una probabilità binomiale con $n$ grande & l'approssimazione gaussiana \vv{tcl} \\
  parla di frequenze relative di un evento o di un simbolo & \vv{frequenzeprob} \\
  chiede in che senso una successione di variabili converge & \vv{convergenze} \\
  \bottomrule
\end{tabular}

Il primo passo è sempre lo stesso: riconoscere che i dati sono i.i.d. (\vv{iid}) e trovare $\mu$ e $\sigma^2$ di un
dato singolo. Il resto viene da lì.
\end{scelta}

\scheda{iid}{
  \usoquando{I dati vengono dallo stesso esperimento ripetuto da capo, senza che un risultato influenzi gli altri:
    lanci, misure ripetute, bit trasmessi.}
  \usocome{La congiunta è il prodotto delle marginali, tutte uguali; basta conoscere la legge di un dato.}
  \usoesempio{Dieci lanci di un dado: $X_1, \ldots, X_{10}$ i.i.d., uniformi su $\{1, \ldots, 6\}$, con $\mu =
    3{,}5$ e $\sigma^2 = \frac{35}{12}$.}
}

\scheda{mediacampione}{
  \usoquando{Si mediano $n$ dati per avvicinarsi alla media vera: misure ripetute, frequenze, sondaggi.}
  \usocome{La media non cambia, la varianza si divide per $n$: la deviazione standard cala come
    $\frac{1}{\sqrt{n}}$, quindi per dimezzarla servono quattro volte più dati.}
  \usoesempio{Cento lanci di un dado: $E\left[ \overline{X}_{100} \right] = 3{,}5$ e $\mathrm{VAR} =
    \frac{35}{12 \cdot 100} \simeq 0{,}029$, cioè deviazione standard $\simeq 0{,}17$.}
}

\scheda{lgn}{
  \usoquando{Si chiede quanti dati servono per garantire, qualunque sia la legge dei dati, che la media campione cada
    vicino alla media vera.}
  \usocome{Si scrive il limite di Chebyshev con $\sigma^2$ di un dato e la tolleranza $\epsilon$, e si impone che sia al
    più la probabilità d'errore $\delta$ ammessa.}
  \usoattenzione{Il limite è una garanzia, spesso pessimista: con il teorema centrale del limite bastano di solito
    molti meno dati.}
  \usoesempio{Moneta onesta, $\sigma^2 = \frac{1}{4}$, $\epsilon = 0{,}05$: $P \leq \frac{0{,}25}{0{,}0025\, n} =
    \frac{100}{n}$, che è al più $0{,}05$ con $n \geq 2000$.}
}

\scheda{frequenzeprob}{
  \usoquando{Si stima una probabilità con la frequenza relativa, o si chiede quanto una frequenza osservata può
    discostarsi dalla probabilità.}
  \usocome{La frequenza è la media campione dell'indicatore dell'evento: media $p$, varianza $\frac{p(1 - p)}{n}$.}
  \usoesempio{La faccia 6 in 600 lanci di un dado: $E\left[ \frac{N}{600} \right] = \frac{1}{6}$, $\mathrm{VAR} =
    \frac{5}{36 \cdot 600} \simeq 2{,}3 \cdot 10^{-4}$, deviazione standard $\simeq 0{,}015$.}
}

\scheda{convergenze}{
  \usoquando{Si chiede in quale senso una successione di variabili aleatorie (medie campione, stimatori) tende a un
    valore.}
  \usocome{In probabilità è un limite di numeri (le probabilità); quasi certa è la probabilità di un limite (che la
    singola realizzazione converga); in media quadratica è il limite dell'errore quadratico medio.}
  \usoesempio{Media campione: $E\left[ \left( \overline{X}_n - \mu \right)^2 \right] = \frac{\sigma^2}{n} \to 0$,
    quindi converge in media quadratica, e per Markov anche in probabilità.}
}

\scheda{sommaiid}{
  \usoquando{Si sommano tante variabili uguali e indipendenti: un totale, un numero di successi, un errore che si
    accumula.}
  \usocome{Media e varianza si moltiplicano per $n$; per usare la gaussiana si standardizza.}
  \usoesempio{Dieci dadi: $\mu_Z = 35$, $\sigma_Z^2 = 10 \cdot \frac{35}{12} \simeq 29{,}2$, deviazione standard
    $\simeq 5{,}4$.}
}

\scheda{tcl}{
  \usoquando{Una grandezza è la somma (o la media) di tanti contributi indipendenti, e servono probabilità
    approssimate: totali, medie campione, binomiali con $n$ grande.}
  \usocome{Si calcolano media e varianza della somma, si standardizza e si usa la funzione $Q$ (\vv{codegauss}).}
  \usoattenzione{Serve la varianza finita: per la Cauchy la media di $n$ dati è ancora una Cauchy, larga come un dato
    solo. Sulle code l'approssimazione è meno precisa.}
  \usoesempio{Moneta onesta, $\epsilon = 0{,}05$: $P \simeq 2\, Q(0{,}1 \sqrt{n}) = 0{,}05$ quando $0{,}1\sqrt{n}
    \simeq 1{,}96$, cioè $n \simeq 385$, contro i 2000 di Chebyshev. Con $n = 50$ lanci il numero di teste è circa
    $\mathcal{N}(25;\ 12{,}5)$.}
}

\scheda{generatrice}{
  \usoquando{Per capire perché la somma di tante variabili diventa gaussiana: è lo strumento della dimostrazione del
    teorema centrale del limite.}
  \usocome{Si calcola la media di $e^{tY}$ con la pmf o con la densità; per una somma di variabili indipendenti si
    moltiplicano le $M$.}
  \usoesempio{$Y = \pm 1$ con probabilità $\frac{1}{2}$: $M_Y(t) = \frac{e^t + e^{-t}}{2}$, e la somma di $n$
    variabili così, indipendenti, ha $M(t) = \left( \frac{e^t + e^{-t}}{2} \right)^n$.}
}
